\section{Pierścień i  ciało(algebra) oraz $\sigma-$ pierścień i $\sigma-$ algebra}
\begin{context}
	X niepusta przestrzeń.
\end{context}

\begin{definition}[Pierścień i ciało zbiorów.]
	
	Mówimy, że rodzina $\mathcal{R}\subseteq \mathcal{P}(\left( X \right) ) $ jest \textbf{pierścieniem zbiorów}, jeżeli
	\begin{enumerate}[label=\alph*)]
		\item \textbf{Zawiera zbiór pusty:} $\O\in \mathcal{R}$,
		\item \textbf{Jest zamknięta na skończone sumy i różnice:}  $A,B \in \mathcal{R} \implies A\cup B, A\setminus B \in \mathcal{R}$
	\end{enumerate}

	Mówimy, że rodzina $\mathcal{R}\subseteq \mathcal{P}(X) $ jest \textbf{ciałem(algebrą) zbiorów} jeśli $\mathcal{R}$ jest  \textbf{pierścieniem zbiorów} oraz $\bm{X \in \mathcal{R}}$. 
\end{definition}
W pierścieniu $\mathcal{R}$ możemy wykonywać operację różnicy symetrycznej $A\Delta  B$, dla dowolnych  $A,B \in \mathcal{R}$. Istotnie
\begin{itemize}
	\item Jeżeli $A,B \in \mathcal{R}$ to $A\Delta B = \overset{\Def}{=} A\setminus B \cup B\setminus A \in \mathcal{R}$.
\end{itemize}
Ponadto możemy także brać przekroje:
 \begin{itemize}
	\item $A\cap B = A \setminus \left( A\setminus B \right) \in \mathcal{R}$.
\end{itemize}
\begin{remark}
	Na to, ażeby rodzina $\mathcal{R}\subseteq \mathcal{P}(X) $ była \textbf{ciałem(algebrą) } potrzeba i wystarcza $A\cup B, A^{c} \in \mathcal{R}$, dla dowolnych $A,B \in \mathcal{R}$ .
\end{remark} 
\begin{proof}
	Ustalmy rodzinę $\mathcal{R} \subseteq \mathcal{P}(X)$. Dla dowolnych $A,B \in \mathcal{R}$, operację różnicy mnogościowej oraz zbiór $X$ możemy wyrazić przy pomocy dopełnień w sposób następujący.
	\begin{itemize}
		\item $A\setminus B = A^{c}\cap B = \overset{\text{de Morgan}}{=} \left( A \cup B^{c} \right) ^{c}$,
	        \item $\O^{c} = X$.
	\end{itemize}
\end{proof} 

	Jeżeli dana rodzina zbiorów $\mathcal{R}$ jest zamknięta na sumy dwóch swoich elementów, to na mocy indukcji, $\bigcup_{i=1} ^{n}A_{i}\in \mathcal{R}$ dla dowolnego $n$ i $A_{i}\in \mathcal{R}$.

	Możemy więc powiedzieć, że \textbf{ciało(algebra)} zbiorów to rodzina zamknięta na wszystkie \textbf{skończone} operacje mnogościowe.

	Należy tu zauważyć dosyć często spotykany błąd:
\begin{remark}[Granica Indukcji: Dlaczego $\forall n$ to nie to samo co $n=\infty$?]
    Niech $\mathcal{R}$ będzie pierścieniem. Z definicji wiemy, że $A \cup B \in \mathcal{R}$.
    Stosując indukcję matematyczną, możemy udowodnić, że dla \textbf{każdego skończonego} $N \in \mathbb{N}$:
    \[
        \text{Jeżeli } A_1, \dots, A_N \in \mathcal{R}, \text{ to } \bigcup_{n=1}^{N} A_n \in \mathcal{R}.
    \]
    To oznacza, że mamy ciąg przynależności:
    \[
        S_1 \in \mathcal{R}, \quad S_2 \in \mathcal{R}, \quad \dots, \quad S_{1000} \in \mathcal{R}, \quad \dots
    \]
    Gdzie $S_N$ to suma częściowa.
    
    \textbf{Tu jest pułapka:} Indukcja zapewnia nas, że każdy wyraz tego ciągu z osobna należy do $\mathcal{R}$. Nie daje nam jednak żadnej gwarancji co do \textbf{granicy} tego ciągu (sumy nieskończonej).
    \[
        \bigcup_{n=1}^{\infty} A_n \overset{?}{\in} \mathcal{R}
    \]
    Symbol $\infty$ nie jest liczbą naturalną, więc indukcja do niego "nie dociera". Aby zagwarantować przynależność sumy nieskończonej, potrzebujemy nowego, silniejszego aksjomatu: definicji \textbf{$\sigma$-pierścienia}.
    
    \textbf{Przykład:} W pierścieniu przedziałów $[a, b)$:
    \begin{itemize}
        \item Suma skończona: $\bigcup_{n=1}^N [\frac{1}{n}, 1) = [\frac{1}{N}, 1) \in \mathcal{R}$. (Jest to przedział domknięto-otwarty).
        \item Suma nieskończona: $\bigcup_{n=1}^\infty [\frac{1}{n}, 1) = (0, 1) \notin \mathcal{R}$. (Wynik "uciekł" z pierścienia, stając się zbiorem otwartym).
    \end{itemize}
\end{remark}
\begin{definition}[$\sigma-$pierścień i $\sigma-$ciało zbiorów.]

	Mówimy, że $\mathcal{R}\subseteq \mathcal{P}(X) $ jest \textbf{$\sigma$-pierścieniem zbiorów} jeżeli $\mathcal{R}$ jest \textbf{zamknięta na przeliczalne sumy}, to znaczy
	\[
		\bigcup_{n=1} ^{\infty}A_n \in \mathcal{R} \text{ dla dowolnego ciągu }A_n \in \mathcal{R}
	.\] 
	
	Jeżeli $\mathcal{R}$ jest $\sigma$ \textbf{-pierścieniem} oraz $X\in \mathcal{R}$ to $\mathcal{R}$ nazywamy $\bm{\sigma-\text{ciałem}}$.
\end{definition}

\begin{remark}[Zamkniętość $\sigma$-ciała na operacje przeliczalne]
    W $\sigma$-ciele $\mathcal{R}$ wykonywane są wszystkie przeliczalne operacje mnogościowe.
    
    Istotnie, niech $\Sigma$ będzie $\sigma$-ciałem. Z definicji wiemy, że jeżeli $A_n \in \Sigma$, to $\bigcup_{n=1}^\infty A_n \in \Sigma$ oraz $A_n^c \in \Sigma$.
    Stąd wynika wykonalność wszystkich innych przeliczalnych operacji:
    
    \begin{enumerate}
        \item \textbf{Przeliczalne przekroje:} Wynikają z praw de Morgana.
        \[
            \bigcap_{n=1}^{\infty} A_n = \left( \bigcup_{n=1}^{\infty} A_n^c \right)^c \in \Sigma.
        \]
        (Ponieważ $A_n^c \in \Sigma \implies \bigcup A_n^c \in \Sigma \implies (\dots)^c \in \Sigma$).
        
        \item \textbf{Granice zbiorów:} Wynikają z definicji konstrukcyjnych, będących kombinacją sum i przekrojów.
        \[
            \limsup_{n \to \infty} A_n \overset{\Def}{=} \bigcap_{n=1}^{\infty}\bigcup_{k=n}^{\infty}A_{n}  \in \Sigma, \quad
            \liminf_{n \to \infty} A_n \overset{\Def}{=}\bigcup_{n=1}^{\infty}\bigcap_{k=n} ^{\infty}  \in \Sigma.
        \]
	Ponieważ $\sigma$-ciało jest zamknięte na przeliczalne sumy i przeliczalne przekroje, a granice zbiorów ($\lim\sup$ oraz $\lim\inf$ są zdefiniowane jako kombinacje tych dwóch operacji, więc granice górna i dolna także należą do $\sigma$-ciała.
    \end{enumerate}
    \textbf{Wniosek:} $\sigma$-ciało jest bezpiecznym środowiskiem dla Analizy – nie można z niego "wypaść" stosując przejścia graniczne na zbiorach.
\end{remark}
\begin{eg}
	Rodzina $\mathcal{R}=\{\O\} $ jest oczywiście pierścieniem, a rodzina $\mathcal{A}=\{\O, X\} $ jest najmniejszym ciałem podzbiorów $X$.
\end{eg} 
\begin{eg}
	Zbiór potęgowy $\mathcal{P}(X) $ jest $\sigma$-ciałem(algebrą).

	Dodatkowo $\mathcal{P}(X) $ to \textbf{maksymalne możliwe $\sigma$-ciało} na przestrzeni $X$.
	\begin{itemize}
		\item Każde inne $\sigma$-ciało $\Sigma$ jest podzbiorem $\mathcal{P}(X) $ ($\Sigma \times \mathcal{P}(X) $ ).
		\item Jest to struktura ,,najbogatsza''- nic więcej do niej nie da się dodać
	\end{itemize}
\end{eg} 
\begin{eg}
	Jeśli oznaczymy przez $\mathcal{R}$ rodzinę wszystkich skończonych podzbiorów nieskończonej przestrzeni $X$ to $\mathcal{R}$ jest pierścieniem, ale nie jest ciałem.

	Istotnie $\O \in \mathcal{R}$ oraz dla dowolnych $A,B \in \mathcal{R}$ $A\cup B$ jest zbiorem skończonym, jak również $A\setminus B$. Zatem $\mathcal{R}$ jest pierścieniem. 
	
	Inaczej, ma się z byciem ciałem. Załóżmy niewprost, że $\mathcal{R}$ jest ciałem, wtedy $X\in \mathcal{R}$. No i tu już jest sprzeczność bowiem $X\subseteq X$, ale $X$ jest nieskończony, więć nie należy do $\mathcal{R}$.

	Zauważmy, że tak zdefiniowne $\mathcal{R}$ nie jest $\sigma$-pierścieniem. Istotnie, załóżmy, że jest. Wtedy w  $X$ możemy wyróżnić różnowarościowy ciąg $x_{n}$. Połóżmy teraz $A_{n}= \{x_{k}: k<n\} $. Każdy element $A_{n}$ jest skończony i co za tym idzie, jest elementem $\mathcal{R}$.
	Jednak $\bigcup_{i \in  I}^{\infty} A_{n} = {x_{n}}_{n=1}^{\infty}$, jest nieskończona. Nie może się zatem zawierać w $\mathcal{R}$.
	
	Analogicznie w nieprzeliczalnej przestrzeni $X$, rodzina $\mathcal{C}$ wszystkich podbiorów przeliczalnych stanowi naturalny przykład $\sigma$- pierścienia, który nie jest $\sigma$-ciałem.
\end{eg} 
\begin{theorem}[Skończona suma przedziałów jest pierścieniem.]
	Niech $\mathcal{I}$ będzie rodziną przedziałów prostej rzeczywistej, która są lewostronnie domknięte:
	\[
	\mathcal{I} = \{[a,b): a,b\in \R\} 
	.\] 
	Niech $\mathcal{R}$ będzie rodziną zbiorów, które są \textbf{skończonymi sumami} elementów z $\mathcal{I}$ :  
	\[
	\mathcal{R} = \left\{\bigcup_{k=1}^{n}I_{k}: n\in \N \text{ oraz }I_{k}\in \mathcal{I}\right\} 
	.\] 
	Wtedy 
	\begin{enumerate}[label=\arabic*.]
		\item $\mathcal{R}$  jest \textbf{pierścieniem} .
		\item Każdy zbiór $A\in \mathcal{R}$ ma takie przedstawienie, w którym przedziały $I_{k}$ są \textbf{parami rozłączne}. 
	\end{enumerate}
\end{theorem}
\vspace{1em}

\hrule

\vspace{1em}

\noindent \textbf{Początek dowodu}

Tendencyjnie, dowód zostanie przeprowadzony nad wyraz formalnie, w celu ujęcia waloru technicznego i- co za tym idzie- edukacyjnego.
Najpierw udowodnimy pomocnicze lematy:

\begin{lemma}[{(i) Różnica dwóch przedziałów jest elementem $\mathcal{R}$}]
	Dla dowolnych przedziałów $I,J \in \mathcal{I}$ ich różnica $I\setminus J$ należy do $\mathcal{R}$.
\end{lemma} 
\begin{proof}
	Niech $I = [a, b)$ oraz $J = [c, d)$. Rozważmy możliwe relacje między tymi przedziałami. Poniższa lista wyczerpuje wszystkie przypadki wzajemnego położenia końców przedziałów.

    \begin{enumerate}[label=\textbf{\arabic*.}]
        
        % PRZYPADEK 1: ROZŁĄCZNE
        \item \textbf{Rozłączne} ($I \cap J = \emptyset$).\\
        Sytuacja, gdy $J$ leży całkowicie na lewo lub na prawo od $I$. Wtedy odjęcie $J$ nie zmienia $I$.
        \[
            I \setminus J = I = [a, b) \in \mathcal{I} \subseteq \mathcal{R}.
        \]
        
        \begin{center}
        \begin{tikzpicture}[xscale=1.5]
            % Oś
            \draw[axis] (-0.5,0) -- (6,0) node[right] {$\R$};
            
            % Przedział I (a, b) = (1, 3)
            \node[dot, blue] (a) at (1, 0.4) {}; 
	    \node[open, blue] (b) at (3, 0.4) {};

	    \draw[intervalI] (a) -- (b) node[midway, above] {$I$};
                       
            % Przedział J (c, d) = (4, 5.5)
            \node[dot, red] (c) at (4, -0.4) {}; 
	    \node[open, red] (d) at (5.5, -0.4) {};
            \draw[intervalJ] (c) -- (d) node[midway, below] {$J$};

            
            % Wynik na osi (a, b)
            \node[dot, green!60!black] (e) at (1, 0) {}; 
	    \node[open, green!60!black] (f) at (3, 0) {};
	    \draw[result] (e) -- (f);
           
            % Opisy
            \node[below] at (1, -0.1) {$a$}; \node[below] at (3, -0.1) {$b$};
            \node[below] at (4, -0.6) {$c$}; \node[below] at (5.5, -0.6) {$d$};
        \end{tikzpicture}
        \end{center}

        % PRZYPADEK 2: ZAWIERANIE
        \item \textbf{Zawieranie} ($I \subseteq J$).\\
        Sytuacja, gdy przedział $J$ całkowicie przykrywa przedział $I$. Wtedy nic nie zostaje.
        \[
            I \setminus J = \emptyset = [0, 0) \in \mathcal{I} \subseteq \mathcal{R}.
        \]

        \begin{center}
        \begin{tikzpicture}[xscale=1.5]
            \draw[axis] (-0.5,0) -- (6,0);
            
            % Przedział I (2, 4)
            \node[dot, blue] (a) at (2, 0.4) {}; \node[open, blue] (b) at (4, 0.4) {};
            \draw[intervalI] (a) -- (b) node[midway, above] {$I$};

            
             Przedział J (1, 5) - przykrywa I
            \node[dot, red] (c) at (1, -0.4) {}; \node[open, red] (d) at (5, -0.4) {};
            \draw[intervalJ] (c) -- (d) node[midway, below] {$J$};

            
            % Linie pomocnicze
            \draw[dashedline] (2, 0.4) -- (2, -0.4);
            \draw[dashedline] (4, 0.4) -- (4, -0.4);
            
            % Wynik: Pusty (brak zielonej linii)
            \node[green!60!black] at (3, 0.15) {$\bm{ \emptyset } $};
            
            \node[below] at (2, 0) {$a$}; \node[below] at (4, 0) {$b$};
            \node[below] at (1, -0.6) {$c$}; \node[below] at (5, -0.6) {$d$};
        \end{tikzpicture}
        \end{center}

        % PRZYPADEK 3: ODCIĘCIE LEWEJ STRONY
        \item \textbf{Odcięcie lewej strony} ($c \le a < d < b$).\\
        Przedział $J$ zachodzi na początek $I$, ucinając go do punktu $d$.
        \[
            [a, b) \setminus [c, d) = [d, b) \in \mathcal{I} \subseteq \mathcal{R}.
        \]

        \begin{center}
        \begin{tikzpicture}[xscale=1.5]
            \draw[axis] (-0.5,0) -- (6,0);
            
            % Przedział I (2, 5)
            \node[dot, blue] (a) at (2, 0.4) {}; \node[open, blue] (b) at (5, 0.4) {};
            \draw[intervalI] (a) -- (b) node[midway, above right] {$I$};

            
            % Przedział J (1, 3.5)
            \node[dot, red] (c) at (1, -0.4) {}; \node[open, red] (d) at (3.5, -0.4) {};
            \draw[intervalJ] (c) -- (d) node[midway, below] {$J$};

            
            % Wynik (3.5, 5) - od d do b
            \node[dot, green!60!black] (e) at (3.5, 0) {}; \node[open, green!60!black] (f) at (5, 0) {};
            \draw[result] (e) -- (f);

            
            \draw[dashedline] (3.5, -0.4) -- (3.5, 0.4);
            
            \node[below] at (2, 0.35) {$a$}; \node[below] at (5, -0.1) {$b$};
            \node[below] at (1, -0.6) {$c$}; \node[below] at (3.5, -0.6) {$d$};
        \end{tikzpicture}
        \end{center}

        % PRZYPADEK 4: ODCIĘCIE PRAWEJ STRONY
        \item \textbf{Odcięcie prawej strony} ($a < c < b \le d$).\\
        Przedział $J$ zachodzi na koniec $I$, ucinając go od punktu $c$.
        \[
            [a, b) \setminus [c, d) = [a, c) \in \mathcal{I} \subseteq \mathcal{R}.
        \]

        \begin{center}
        \begin{tikzpicture}[xscale=1.5]
            \draw[axis] (-0.5,0) -- (6,0);
            
            % Przedział I (1, 4)
            \node[dot, blue] (a) at (1, 0.4) {}; \node[open, blue] (b) at (4, 0.4) {};
            \draw[intervalI] (a) -- (b) node[midway, above left] {$I$};

            
            % Przedział J (3, 5)
            \node[dot, red] (c) at (3, -0.4) {}; \node[open, red] (d) at (5, -0.4) {};
            \draw[intervalJ] (c) -- (d) node[midway, below] {$J$};

            % Wynik (1, 3) - od a do c
            \node[dot, green!60!black] (e) at (1, 0) {}; \node[open, green!60!black] (f) at (3, 0) {};
            \draw[result] (e) -- (f);

            
            \draw[dashedline] (3, -0.4) -- (3, 0.4);
            
            \node[below] at (1, -0.1) {$a$}; \node[below] at (4, 0.075) {$b$};
            \node[below] at (3, -0.6) {$c$}; \node[below] at (5, -0.6) {$d$};
        \end{tikzpicture}
        \end{center}

        % PRZYPADEK 5: WYCIĘCIE ŚRODKA
        \item \textbf{Wycięcie środka} ($a < c < d < b$).\\
        Przedział $J$ zawiera się ściśle wewnątrz $I$. W wyniku otrzymujemy sumę dwóch rozłącznych przedziałów lewostronnie domkniętych.
        \[
            [a, b) \setminus [c, d) = [a, c) \cup [d, b).
        \]
        Jest to suma dwóch elementów z $\mathcal{I}$, a zatem z definicji pierścienia generowanego, wynik należy do $\mathcal{R}$.

        \begin{center}
        \begin{tikzpicture}[xscale=1.5]
            \draw[axis] (-0.5,0) -- (6,0);
            
            % Przedział I (1, 5.5)
            \node[dot, blue] (a) at (1, 0.4) {}; \node[open, blue] (b) at (5.5, 0.4) {};
            \draw[intervalI] (a) -- (b) node[pos=0.1, above] {$I$};

            
            % Przedział J (2.5, 4) - w środku
            \node[dot, red] (c) at (2.5, -0.4) {}; \node[open, red] (d) at (4, -0.4) {};
            \draw[intervalJ] (c) -- (d) node[midway, below] {$J$};
            
            % Wynik: [a,c) suma [d,b)
            % Lewa część
            \node[dot, green!60!black] (e) at (1, 0) {}; \node[open, green!60!black] (f) at (2.5, 0) {};
            \draw[result] (e) -- (f);
            
            % Prawa część
            \node[dot, green!60!black] (g) at (4, 0) {}; \node[open, green!60!black] (h) at (5.5, 0) {};
            \draw[result] (g) -- (h);
            
            % Linie pomocnicze
            \draw[dashedline] (2.5, -0.4) -- (2.5, 0.4);
            \draw[dashedline] (4, -0.4) -- (4, 0.4);
            
            \node[below] at (1, -0.1) {$a$}; \node[below] at (5.5, -0.1) {$b$};
            \node[below] at (2.5, -0.6) {$c$}; \node[below] at (4, -0.6) {$d$};
        \end{tikzpicture}
        \end{center}

    \end{enumerate}
    
    W każdym przypadku wynik operacji jest skończoną sumą przedziałów z $\mathcal{I}$, co kończy dowód.
\end{proof} 

\begin{lemma}[{(ii) Zamkniętość $\mathcal{R}$ na sumy skończone}]
    Niech $A_1, \dots, A_n$ będą zbiorami należącymi do rodziny $\mathcal{R}$. Wtedy ich suma $\bigcup_{i=1}^{n} A_i$ również należy do $\mathcal{R}$.
\end{lemma} 
\begin{proof}
	    Dowód przeprowadzimy przez indukcję względem liczby zbiorów $n$.

    \textbf{1. Baza indukcji ($n=1$):} \\
    Oczywiste jest, że $\bigcup_{i=1}^{1} A_i = A_1 \in \mathcal{R}$.

    \textbf{2. Krok pomocniczy (dla dwóch zbiorów):} \\
    Wykażemy najpierw, że suma dwóch zbiorów z $\mathcal{R}$ należy do $\mathcal{R}$.
    Niech $A, B \in \mathcal{R}$. Z definicji rodziny $\mathcal{R}$, zbiory te są skończonymi sumami przedziałów (generatorów z $\mathcal{I}$). Istnieją zatem liczby $p, q \in \mathbb{N}$ oraz przedziały $I_1, \dots, I_p \in \mathcal{I}$ i $J_1, \dots, J_q \in \mathcal{I}$ takie, że:
    \[
        A = \bigcup_{k=1}^{p} I_k \quad \text{oraz} \quad B = \bigcup_{k=1}^{q} J_k.
    \]
    Zdefiniujmy nowy ciąg przedziałów $(K_m)_{m=1}^{p+q}$, dokonując reindeksacji (sklejenia obu rodzin):
    \[
        K_m = 
        \begin{cases}
            I_m & \text{dla } 1 \le m \le p, \\
            J_{m-p} & \text{dla } p < m \le p+q.
        \end{cases}
    \]
    Zauważmy, że suma tego nowego ciągu to dokładnie $A \cup B$:
    \[
        \bigcup_{m=1}^{p+q} K_m = \left( \bigcup_{k=1}^p I_k \right) \cup \left( \bigcup_{k=1}^q J_k \right) = A \cup B.
    \]
    Ponieważ każdy $K_m$ jest przedziałem (elementem $\mathcal{I}$), to $A \cup B$ jest skończoną sumą przedziałów, co oznacza, że $A \cup B \in \mathcal{R}$.

    \textbf{3. Krok indukcyjny ($n \to n+1$):} \\
    Załóżmy, że twierdzenie jest prawdziwe dla dowolnego układu $n$ zbiorów z $\mathcal{R}$. Rozważmy $n+1$ zbiorów $A_1, \dots, A_{n+1} \in \mathcal{R}$.
    Zapiszmy ich sumę jako:
    \[
        \bigcup_{i=1}^{n+1} A_i = \left( \bigcup_{i=1}^{n} A_i \right) \cup A_{n+1}.
    \]
    Oznaczmy $S_n = \bigcup_{i=1}^{n} A_i$. Z założenia indukcyjnego wiemy, że $S_n \in \mathcal{R}$.
    Mamy więc sumę dwóch elementów pierścienia: $S_n$ oraz $A_{n+1}$.
    Na mocy udowodnionego w punkcie 2 kroku pomocniczego, suma ta należy do $\mathcal{R}$.
    
    Na mocy zasady indukcji matematycznej twierdzenie jest prawdziwe dla każdego $n$.
\end{proof} 

\begin{lemma}[{(iii) Różnica przedziału z $\mathcal{I}$ oraz zbioru z $\mathcal{R}$}]
	Dla dowolnego przedziału $I\in \mathcal{I}$ oraz dowolnego zbioru $B\in \mathcal{R}$ różnica $I\setminus B$ należy do $\mathcal{R}$.
\end{lemma} 
\begin{proof}
	Ustalmy dowolny zbiór $I\in \mathcal{I}$. Przeprowadzimy dowód przez indukcję względem liczby elementów zbioru $B$. Niech $B = \bigcup_{j=1} ^{m}J_{j}$ gdzie $J_{j}\in \mathcal{I}$.

\textbf{Baza indukcji} (m=1): Niech $B= J_{j}$. Wtedy $I\setminus B = I\setminus J_1$. Na mocy Lematu (i) $I\setminus J_1 \in \mathcal{R}$. Zatem baza jest spełniona.

\textbf{Założenie:} Załóżmy, że dla dowolnej sumy $m$ przedziałów $B_{m}$ różnica $I\setminus B_{m}\in \mathcal{R}$.

\textbf{Krok indukcyjny ($m\to m+1$):} Niech $B_{m+1}= \bigcup_{j=1} ^{m+1}J_{j} = \left( \bigcup_{j=1} ^{m} J_{j}\right) \cup J_{m+1} $. Wtedy
\[
I\setminus B_{m+1} = I\setminus \left( B_{m}\cup J_{m+1} \right) \overset{\text{tożsamość}}{=} \left( I\setminus B_{m} \right) \setminus J_{m+1}
.\] 
    
Z założenia indukcyjnego, $I\setminus B_{m}\in \mathcal{R}$. Oznacza to, że istnieje pewne $k\in N$ oraz przedziały  $K_1,K_2,\ldots,K_{k}\in \mathcal{R}$, takie że
\[
I\setminus B_{m} = \bigcup_{p=1}^{k}K_{p}
.\] 
Podstawiając to do wzoru wyjściowego:
\[
	\left( I\setminus B_{m} \right) \setminus J_{m+1} = \left( \bigcup_{p=1} ^{k}K_{p} \right) \setminus J_{m+1}
.\] 
Korzystając z prawa rozdzielnośći różnicy względem sumy ($\bigcup_{i}  A_{i} \setminus C = \bigcup_{i}  \left( A_{i} \setminus C\right)$):
\[
\bigcup_{p=1} ^{k}K_{p}\setminus J_{m+1} = \bigcup_{p=1} ^{k}\left( K_{p}\setminus J_{m+1} \right) 
.\] 

Rozważmy teraz pojedynczy składnik sumy $\bigcup_{p=1} ^{k}\left( K_{p}\setminus J_{m+1} \right) $, mianowicie $K_{p}\setminus J_{m+1}$. $K_{p}$ jest przedziałem oraz $J_{m+1}$ jest przedziałem, zatem na mocy Lematu (i) jest on elementem $\mathcal{R}$. 

Powyższa suma jest skończoną($k$ elementów) sumą elementów z $\mathcal{R}$. Zatem na mocy Lematu (ii), owa suma także jest elementem $\mathcal{R}$. Zatem $I\setminus B_{m+1} \in \mathcal{R}$.

Na mocy zasady indukcji matematycznej, twierdzenie jest prawdziwe dla każdego $m$.
\end{proof} 

\vspace{1em}

\noindent Pokażemy teraz, że $\mathcal{R}$ jest pierścieniem.
\begin{proof}
	Sprawdzamy trzy warunki definiujące pierścień:
	\begin{enumerate}[label=\alph*)]
		\item Dla dowolnego $a\in \R$, $\O = [a,a) \in \mathcal{I}\subseteq \mathcal{R}$.
	        \item Dla dowolnych $A,B \in \mathcal{R}$ ich suma $A\cup B$ na mocy Lematu (ii) jest elementem $\mathcal{R}$.
		\item Dla dowolnych $A,B \in \mathcal{R}$ ich różnica $A\setminus B$ jest elementem $\mathcal{R}$:
	Z definicji rodziny $\mathcal{R}$ zbiór $A$ jest skończoną sumą przedziałów:
\[
A = \bigcup_{i=1} ^{n}I_{n}\text{, gdzie } I_{k} \in \mathcal{I}
.\] 

Obliczamy różnicę $A\setminus B$:
\[
	A\setminus B = \left(\bigcup_{i=1} ^{n} I_{k}  \right) \setminus B = \bigcup_{i=1} ^{n}\left( I_{k}\setminus B \right)
.\] 

Dla każdego $i\in \{1,\ldots,n\} $ zbiór $I_{i}$ jest przedziałem, a  $B$ jest elementem $\mathcal{R}$. Z Lematu (iii) wnioskujemy, że każdy składnik $I_{i}\setminus B$ jest elementem $\mathcal{R}$.
Ostatecznie zbiór $A\setminus B$ jest skończoną sumą elementów z $\mathcal{R}$. Na mocy Lematu (ii)), wnioskujemy, że $A\setminus B \in \mathcal{R}.$
	\end{enumerate} 
\end{proof} 

\vspace{1em}

\noindent Teraz udowodnimy, że każdy zbiór $A\in \mathcal{R}$ ma takie przedstawienie, w którym przedziały $A_{k}$ są \textbf{rozłączne}.
\begin{proof}
	Niech $A\in \mathcal{R}$, czyli $A = \bigcup_{i=1} ^{n}A_{i}$, gdzie $A_{i} \in  \mathcal{I}$.
	Zdefiniujmy następujący układ zbiorów $B_{i}$ :
	\begin{itemize}
		\item $B_1 = A_1$.
		\item $B_2 = A_2 \setminus A_1$.
		\item $B_3 = A_3\setminus \left( A_2 \cup A_1 \right) $.
		\item $B_4 = A_4\setminus \left( A_3\cup A_2 \cup A_1 \right) $.
		\item $\vdots$
		\item  $B_{n} = A_{n}\setminus \left( \bigcup_{k=1} ^{n-1}A_{k} \right) $
	\end{itemize}

    Niech $S_A = \bigcup_{n=1}^{\infty} A_n$ (suma ,,brudna'') oraz $S_B = \bigcup_{n=1}^{\infty} B_n$ (suma ,,czysta''), gdzie zbiory $B_n$ zdefiniowane są wzorem:
    \[
        B_1 = A_1, \quad B_n = A_n \setminus \bigcup_{i=1}^{n-1} A_i \quad \text{dla } n > 1.
    \]

    \noindent\textbf{1. W jedną stronę ($S_B \subseteq S_A$):} \\
    Z definicji zbioru $B_n$ mamy:
    \[
        B_n = A_n \setminus \left( \bigcup_{i=1}^{n-1} A_i \right).
    \]
    Wynika stąd bezpośrednio, że $B_n \subseteq A_n$ dla każdego $n \in \mathbb{N}$.
    Ponieważ każdy składnik sumy $S_B$ jest podzbiorem odpowiadającego mu składnika sumy $S_A$, to suma podzbiorów musi zawierać się w sumie nadzbiorów:
    \[
        \bigcup_{n=1}^{\infty} B_n \subseteq \bigcup_{n=1}^{\infty} A_n \implies S_B \subseteq S_A.
    \]

    \noindent\textbf{2. W drugą stronę ($S_A \subseteq S_B$):} \\
    Weźmy dowolny punkt $x \in S_A$. Oznacza to, że $x$ należy do przynajmniej jednego zbioru z ciągu $(A_n)$.
    Zdefiniujmy zbiór indeksów $K = \{ n \in \mathbb{N} : x \in A_n \}$. Skoro $x \in S_A$, to $K \neq \emptyset$.
    
    Z zasady minimum dla liczb naturalnych w zbiorze $K$ istnieje element najmniejszy. Oznaczmy go przez $k = \min\{K\} $.
    
    Co to oznacza?
    \begin{itemize}
        \item $x \in A_k$ (ponieważ $k \in K$),
        \item $x \notin A_1, x \notin A_2, \dots, x \notin A_{k-1}$ (ponieważ $k$ jest najmniejszym indeksem, dla którego $x \in A_k$).
    \end{itemize}
    Zatem $x$ nie należy do sumy wcześniejszych zbiorów:
    \[
        x \notin \bigcup_{i=1}^{k-1} A_i.
    \]
    Łącząc te dwa fakty ($x \in A_k$ i $x \notin \bigcup_{i=1}^{k-1} A_i$), otrzymujemy z definicji zbioru $B_k$:
    \[
        x \in A_k \setminus \bigcup_{i=1}^{k-1} A_i \implies x \in B_k.
    \]
    Skoro $x \in B_k$, to tym bardziej $x \in \bigcup_{n=1}^{\infty} B_n$, czyli $x \in S_B$.
    
    \noindent\textbf{Wniosek:} \\
    Każdy punkt z $S_A$ trafia do $S_B$ (konkretnie do swojej ,,pierwszej'' szufladki $B_k$). Żaden punkt nie ginie, żaden nowy nie powstaje. Zatem:
    \[
        \bigcup_{n=1}^{\infty} A_n = \bigcup_{n=1}^{\infty} B_n.
    \]

    \noindent\textbf{3. Uzasadnienie rozłączności ($B_n$ są parami rozłączne):} \\
    Weźmy dwa różne indeksy $n, m \in \mathbb{N}$ i załóżmy bez straty ogólności, że $m < n$.
    Sprawdźmy przekrój $B_n \cap B_m$.
    Z definicji $B_n$:
    \[
        B_n = A_n \setminus \left( A_1 \cup \dots \cup A_m \cup \dots \cup A_{n-1} \right).
    \]
    Zauważmy, że zbiór $A_m$ jest jednym ze składników sumy, którą odejmujemy. Zatem $B_n$ jest z definicji rozłączny z $A_m$:
    \[
        B_n \cap A_m = \emptyset.
    \]
    Wiemy też z części 1, że $B_m \subseteq A_m$. Skoro $B_n$ jest rozłączny z całym $A_m$, to musi być też rozłączny z jego podzbiorem $B_m$.
    \[
        B_n \cap B_m \subseteq B_n \cap A_m = \emptyset \implies B_n \cap B_m = \emptyset.
    \]
    Zbiory $B_n$ są zatem parami rozłączne.
\end{proof} 
\vspace{1em}

\noindent \textbf{Koniec dowodu.}

\vspace{1em}

\hrule

\vspace{1em}

\noindent Opisanie wprost (konstrukcyjnie) rodziny zbiorów zamkniętej na przeliczalne sumy bywa trudne. Dlatego wygodniej jest definiować takie rodziny poprzez ich \textbf{generatory}.

Aby zastosować powyższy mechanizm do definicji (1.1.5 PlbSkrypt), musimy udowodnić, że struktury algebraiczne(pierścienie, $\sigma$-pierścienie, ciała, $\sigma$-pierścienie) są niezmiennicze na przekroje.

\begin{theorem}
    Niech $\{\mathcal{S}_i\}_{i \in I}$ będzie dowolną rodziną $\sigma$-pierścieni na przestrzeni $X$. Wtedy ich przekrój $\mathcal{S} = \bigcap_{i \in I} \mathcal{S}_i$ również jest $\sigma$-pierścieniem.
\end{theorem}

\begin{proof}
	Chcemy dowieść stabilności struktury oraz jej minimalności.

	Inaczej, chcemy pokazać definicja jest poprawna (definiuje obiekt istniejący, unikalny i najmniejszy).
\end{proof}

Dla dowolnej rodziny $\mathcal{F} \subseteq \mathcal{P}(X)$ definiujemy:
    \begin{itemize}
        \item $\bm{r(\mathcal{F})}$  najmniejszy \textbf{pierścień} zawierający $\mathcal{F}$ (\textbf{R}ing):
		\[
		r\left( \mathcal{F} \right) := \bigcap \{\mathcal{R}: \mathcal{R}\text{ jest pierścieniem oraz }\mathcal{F}\subseteq \mathcal{R}\}  
		.\] 
        \item $\bm{s(\mathcal{F})}$ najmniejszy \textbf{$\sigma$-pierścień} zawierający $\mathcal{F}$ (\textbf{S}igma ring):
		\[
			s\left( \mathcal{F} \right) := \bigcap \{\mathcal{R}: \mathcal{R}\text{ jest $\sigma$-pierścieniem oraz }\mathcal{F}\subseteq \mathcal{R}\} 
		.\] 
        \item $\bm{a(\mathcal{F})}$  najmniejsze \textbf{ciało} (algebrę) zawierające $\mathcal{F}$ (\textbf{A}lgebra):
		\[
		a\left( \mathcal{F} \right) := \bigcap \{\Sigma : \Sigma \text{ jest ciałem oraz }\mathcal{F}\subseteq \Sigma \}
		.\] 
	\item $\bm{\sigma(\mathcal{F})}$  najmniejsze \textbf{$\sigma$-ciało} ($\sigma$-algebra) zawierające $\mathcal{F}$ ($\bm{\sigma}$-algebra):
		\[
		\sigma\left( \mathcal{F} \right) := \bigcap \{\Sigma : \Sigma \text{ jest $\sigma$-ciałem oraz }\mathcal{F}\subseteq \Sigma \}
		.\] 
    \end{itemize}

    \begin{definition}[$\sigma$-ciało zbiorów borelowskich]
    Niech $\mathcal{U}$ oznacza rodzinę wszystkich otwartych podzbiorów prostej $\mathbb{R}$.
    \textbf{$\sigma$-ciałem zbiorów borelowskich} nazywamy najmniejsze $\sigma$-ciało zawierające rodzinę $\mathcal{U}$. Oznaczamy je symbolem $Bor(\mathbb{R})$:
    \[
        Bor(\mathbb{R}) = \sigma(\mathcal{U}).
    \]
\end{definition}
\begin{lemma}[Monotoniczność operacji generowania]
    Niech $X$ będzie przestrzenią oraz $G_1, G_2 \subseteq \mathcal{P}(X)$. 
    Jeżeli $G_1 \subseteq \sigma(G_2)$, to $\sigma(G_1) \subseteq \sigma(G_2)$.
\end{lemma}

\begin{proof}

    TODO

    Zdefiniujmy rodzinę wszystkich $\sigma$-ciał zawierających zbiór $G_1$:
    \[
        \mathcal{S} = \{ \Sigma \subseteq \mathcal{P}(X) : \Sigma \text{ jest } \sigma\text{-ciałem} \land G_1 \subseteq \Sigma \}.
    \]
    Z definicji, $\sigma(G_1) = \bigcap \mathcal{S}$.
    
    Zauważmy teraz dwa fakty dotyczące zbioru $\sigma(G_2)$:
    \begin{enumerate}
        \item $\sigma(G_2)$ jest $\sigma$-ciałem (z definicji operacji generowania).
        \item $G_1 \subseteq \sigma(G_2)$ (z założenia twierdzenia).
    \end{enumerate}
    
    Powyższe warunki oznaczają, że $\sigma(G_2)$ jest elementem rodziny $\mathcal{S}$ ($\sigma(G_2) \in \mathcal{S}$).
    Ponieważ przekrój rodziny zbiorów zawiera się w każdym jej elemencie, otrzymujemy tezę:
    \[
        \sigma(G_1) = \bigcap \mathcal{S} \subseteq \sigma(G_2).
    \]
\end{proof}
\begin{lemma}
	Niech $\mathcal{F}$ będzie rodziną przedziałów $[p,q)$ gdzie $p,q\in \Q$. Wtedy $\sigma\left( \mathcal{F} \right) = \bor\left(\R\right) $
\end{lemma} 
\begin{proof}
	Pokażemy zawieranie w obie strony.

	\noindent$\bm{ \sigma\left( \mathcal{F} \right)  \subseteq \bor\left(\R\right) }$ :

	Ustalmy dowolny przedział $[a,b)\in \mathcal{F}$. Zauważmy, że $\left( a,b \right) = \bigcap_{n=1} ^{\infty}\left( a - \frac{1}{n},b \right)$. Przeliczalny przekrój zbiorów przedziałów otwartych z definicji jest elementem $\bor\left(\R\right) $. Zatem \[
	\mathcal{F}\subseteq \bor\left(\R\right) 
	.\] 
	Na mocy monotoniczności operacji generowania $\sigma\left( \mathcal{F} \right) \subseteq \bor\left(\R\right) $.

	\noindent$\bm{ \bor\left(\R\right)  \subseteq \sigma\left( \mathcal{F} \right):}$ 

	Zauważmy, że z definicji $\bor\left(\R\right) = \sigma\left( \mathcal{I} \right) $, gdzie $\mathcal{I}$ to 
	\[
	\mathcal{I} = \{\left( a,b \right) : a,b \in \R\} 
	.\] 
	Ustalmy dowolny przedział $\left( a,b \right) \in \R$. Wiemy, że dowolny przedział otwarty (zbiór) możemy rozłożyć na przeliczalną sumą przedziałów o końcach wymiernych:
	\[
		\left( a,b \right) = \bigcup_{n=1} ^{\infty}\left( a_{n},b_{n} \right) \text{ gdzie } a_{n},b_{n}\in \Q
	.\] 

	Ponadto zauważmy, że dowolny przedział $\left( a_{n},b_{n} \right) $ można zapisać w postaci:
	\[
		\left( a_{n},b_{n} \right) = \bigcup_{k=1} ^{\infty} \left[a_{n}+\tfrac{1}{k}, b_{n}\right)
	.\] 
	Podsumowując, nasz wyjściowy przedział $\left( a,b \right) $ możemy zapisać w postaci:
	\[
		\left( a,b \right) = \bigcup_{n=1}^{\infty}\left( \bigcup_{k=1} ^{\infty} \left[a_{n}+ \tfrac{1}{k}, b_{n}\right) \right)  
	.\] 
	Następnie, korzystając z zamkniętości $\sigma\left( \mathcal{F} \right) $ na przeliczalne sumy, zauważmy, że skoro $[a_{n}+ \tfrac{1}{k}, b_{n})\in \mathcal{F}$, to $\bigcup_{k=1} ^{\infty} [a_{n}+ \tfrac{1}{k}, b_{n}) \in \sigma\left( \mathcal{F} \right).$ 

\noindent Idąc dalej, ponieważ $\bigcup_{k=1} ^{\infty} [a_{n}+ \tfrac{1}{k}, b_{n}) \in \sigma\left( \mathcal{F} \right),$ więc $\bigcup_{n=1}^{\infty}( \bigcup_{k=1} ^{\infty} [a_{n}+ \tfrac{1}{k}, b_{n})) \in \sigma\left( \mathcal{F} \right) .$ 

\noindent Zatem 
\[
		\left( a,b \right) = \bigcup_{n=1}^{\infty}\left( \bigcup_{k=1} ^{\infty} \left[a_{n}+ \tfrac{1}{k}, b_{n}\right) \right)  \in \sigma\left( \mathcal{F} \right) 
.\] 
Zatem na mocy monotoniczności generowania $\sigma\left( \bor\left(\R\right)  \right) \subseteq \sigma\left( \mathcal{F} \right) $.
\end{proof} 
\begin{remark}[Kluczowe rodziny generujące $Bor(\mathbb{R})$]
    Warto pamiętać, że to samo $\sigma$-ciało jest generowane przez inne, prostsze rodziny zbiorów, np.:
    \begin{enumerate}
        \item \textbf{Przedziały otwarte:} Rodzina $\{(a,b) : a,b \in \mathbb{R}\}$.
        \item \textbf{Przedziały domknięte:} Rodzina $\{[a,b] : a,b \in \mathbb{R}\}$.
        \item \textbf{Półproste (Ważne w Probabilistyce!):} 
        Rodzina $\{(-\infty, a] : a \in \mathbb{R}\}$ lub $\{(a, \infty) : a \in \mathbb{R}\}$.
        \item \textbf{Przedziały "techniczne" (do definicji miary):} 
        Rodzina przedziałów lewostronnie domkniętych $\{[a, b) : a,b \in \mathbb{R}\}$.
        \item \textbf{Generatory przeliczalne (Topologiczne):}
        Rodzina przedziałów o końcach \textbf{wymiernych} $\{(p, q) : p,q \in \mathbb{Q}\}$ (lub $[p, q)$).
    \end{enumerate} 
    Dzięki temu, aby sprawdzić mierzalność funkcji $f: X \to \mathbb{R}$, wystarczy sprawdzić przeciwobrazy zbiorów z jednej z tych prostszych rodzin (np. półprostych).
\end{remark}

\subsection{Wielki kompromis: $\bor\left(\R\right) $ jako Arena Mierzalności}
Zasygnalizujemy tu pojecie, które dokładniej zostanie omówione w następnym rozdziale.

\begin{remark}
	\textbf{Problem: Chcemy umieć mierzyć ,,rozmiar'' zbiorów.}

	W  \R(na prostej) mamy intuicyjne pojęcie "rozmiaru", nazywamy je \textbf{długością.} Potrafimy zmierzyć długość przedziału $[a,b]$- jest to b-a.

	Co dalej?
 	\begin{itemize}
 		\item Jaka jest długość zbioru $\{1,2,3\} $? Intuicja mówi 0.
		\item Jaka jest długość dwóch rozłącznych przedziałów $[0,1] \cup [2,3]$? Oczywiście 1 + 1 = 2.
		\item Jak jest długość $\Q\cap [0,1]$? Jest to zbiór przeliczalny. Nasza intuicja podpowiada, że skoro trzy punkty mają długość 0, to przeliczalna suma punktów też powinna mieć długość 0. To wymaga  \textbf{przeliczalnej(sigma) addytywności}:
			\[
			m\left( \bigcup_{n=1}^{\infty} A_{n} \right) = \sum_{n=1}^{\infty} m\left( A_{n} \right)\text{ dla rozłącznych } A_{n}
			.\]
 	\end{itemize}

	Chcemy aby nasza "funkcja miary" miała dwie własności:
	\begin{enumerate}[label=\alph*)]
		\item Potrafiła mierzyć przedziały.
		\item Była $\sigma$-addytywna.
	\end{enumerate}

	\textbf{Komplikacja: Na jakich zbiorach ta funkcja ma działać?}  

	Chcielibyśmy aby działała na wszystkich podzbiorach \R. Niestety przy założeniu Aksjomatu Wyboru, można udowodnić(konstrukcja Vitalego), że nie da się zdefiniować takiej funkcji miary na wszystkich podzbiorach \R, która mierzyłaby przedziały tak jak chcemy i byłaby $\sigma$-addytywna(i do tego niezmiennicza względem przesunięć).

	Musimy iść na kompromis- nie będziemy mierzyć wszystkich zbiorów. Będziemy mierzyć wszystkie "wystarczająco porządne" zbiory. 

	To jest intuicja stojąca za $\bor\left(\R\right) $.

	,,Porządne zbiory''(czyli te które będziemy mierzyć) powinny tworzyć rodzinę, która jest \textbf{zamknięta} na operacje, które chcemy wykonywać.

	Jeśli potrafię zmierzyć zbiory A i B, to chcę też umieć zmierzyć $A\cup B$ oraz $A\cap B$. Jeśli potrafię zmierzyć zbior $A$ to chcę też umieć zmierzyć jego dopełninie- $A^{c}$. A ponieważ nasza miara ma być  $\sigma$- addytywna, to naturalnym też jest wymagać, że jeśli potrafię zmierzyć przeliczalnie wiele zbiorów $A_1,A_2,\ldots$ to potrafię zmierzyć ich sumę $\bigcup_{n=1} ^{\infty}A_{n}$.

	Rodzina zbiorów, która spełnia te trzy warunki(zamkniętość na dopełnienia, przeliczalne sumy i przeliczalne przekroje- ten ostatni warunek wynika z dwóch pierwszych), to jest właśnie $\sigma$-ciało(lub też $\sigma$ algebra).

	To jest nasz ,,\textbf{klub porządnych zbiorów}''.

	\textbf{Wreszcie: Czym jest Bor(\R)?}

	Mamy już maszynerię($\sigma$-ciało). Od czego zacząć? Jakie zbiory na pewno muszą być w naszym "klubie porządnych zbiorów"? Oczywiście \textbf{przedziały}, to od nich zaczęliśmy.

Pomyślmy szerzej. W analize fundamentalną rolę odgrywają \textbf{zbiory otwarte}. Każdy zbiór otwarty w \R, jest \textbf{przeliczalną sumą rozłącznych przedziałów otwartych.} Skoro chcemy mierzyć przedziały i być zamknięci na przeliczalne sumy, to musimy umieć mierzyć zbiory otwarte.

I tu dochodzimy do definicji.

\textbf{Intuicja } $\bor\left(\R\right) $ : Zbiory borelowskie $\bor\left(\R\right) $ to jest \textbf{najmniejszy możliwy ,,klub''($\sigma$- algebra) ,,porządnych zbiorów'', który zawiera wszystkie zbiory otwarte w \R}.

,,Najmniejszy'' oznacza, że nie wrzucamy tam niczego na siłę. Bierzemy zbiory otwarte i dodajemy tam tylko te zbiory, które są absolutnie niezbędne, aby ten "klub" był $\sigma-$ algebrą.

\textbf{Co z tego wynika?}
\begin{itemize}
	\item Skoro $\bor\left(\R\right) $ zawiera \textbf{zbiory otwarte} to musi zawierać ich dopełnienia, czyli \textbf{zbiory domknięte}.
	\item Skoro zawiera \textbf{zbiory domknięte}, to musi zawierać ich \textbf{przeliczalne sumy}(zbiory typu $F_{\sigma}$ ).  
	\item Skoro zawiera \textbf{zbiory otwarte}, to musi zawierać ich \textbf{przeliczalne przekroje}(zbiory typu $G_{\delta}$.
	\item I tak dalej\ldots. Można to kontynuować w nieskończoność(a nawet dalej, używając indukcję pozaskończoną), tworząc coraz bardziej skomplikowane zbiory, które też są "porządne", bo da się je skonstruować ze zbiorów otwartych w przeliczalnej liczbie kroków.
\end{itemize}

\textbf{Podsumowanie:}

Zbiory borelowskie to nasza arena. To jest rodzina podzbiorów \R, która jest wystarczająco bogata, by zawierać \textbf{wszystkie zbiory}, jakie spotkamy w normalnej analizie(otwarte, domknięte, przedziały, punkty, zbiory przeliczalne), a jednocześnie wystarczająco ,,porządne'' aby dało się na nich \textbf{zawsze} określić miarę(jak miara Lebesgue'a)
\end{remark} 


